\begin{theorem}[Strong Induction on a Peano System]
\label{thm:strong-induction-on-peano-system}
Let $(P,S,1)$ be a Peano system, let $<$ be a binary relation on $P$, and let
$\varphi$ be a predicate on $P$. Suppose that $<$ is compatible with the
successor structure in the following sense:
\[
\forall k\in P,\qquad \neg(k<1),
\]
and
\[
\forall k,n\in P,\qquad
k<S(n)\Longleftrightarrow (k<n\lor k=n).
\]
If
\[
\forall n \in P,\qquad
\left[
\left(\forall k \in P,\ k < n \Longrightarrow \varphi(k)\right)
\Longrightarrow \varphi(n)
\right]
\]
then
\[
\forall n \in P,\qquad \varphi(n).
\]
\hyperref[prf:strong-induction-on-peano-system]{Go to proof.}
\end{theorem}
\end{theorembox}
\LeanFormalizes{thm:strong-induction-on-peano-system}{lra-lean}{LRA.VolumeII.PeanoSystems.Induction.Core}{StrongInductionOnPeanoSystem}{pending}

\begin{remark*}[Standard quantified statement]
For a Peano system $(P,S,1)$, a binary relation $<$ on $P$, and a predicate
$\varphi$ on $P$,
\[
\Bigl[
\bigl(\forall k\in P,\ \neg(k<1)\bigr)
\land
\bigl(\forall k,n\in P,\ k<S(n)\Longleftrightarrow(k<n\lor k=n)\bigr)
\land\\
\qquad
\forall n \in P,\quad
\left(
\left(\forall k \in P,\ k < n \Longrightarrow \varphi(k)\right)
\Longrightarrow \varphi(n)
\right)
\Bigr]
\Longrightarrow
\forall n \in P,\ \varphi(n).
\]
\end{remark*}

\begin{remark*}[Predicate reading]
\[
\operatorname{StrongInductionPrinciple}(P,<,\varphi).
\]
\end{remark*}

\begin{remark*}[Interpretation]
Strong induction grants the inductive step access to the entire history below
$n$, not just the single predecessor. It is the proof form behind least
counterexample arguments.
\end{remark*}

\begin{dependencies}
\begin{itemize}
  \item \hyperref[thm:induction-principle-for-peano-system]{Induction Principle for a Peano System}
\end{itemize}
\end{dependencies}
\end{topicbox}

\begin{topicbox}{Inductive Subsets and Predecessors}
\begin{exposition}
Induction is expressed in terms of subsets of a Peano system. Two structural
notions appear repeatedly in that setting: closure under successor and
membership generated from the distinguished first element. Predecessor language
is also convenient once it is known that every non-initial element arises as a
successor.
\end{exposition}

